% Created 2026-05-11 Mon 02:21
% Intended LaTeX compiler: lualatex
\documentclass[11pt]{article}

\usepackage{amsmath}
\usepackage{fontspec}
\usepackage{graphicx}
\usepackage{longtable}
\usepackage{wrapfig}
\usepackage{rotating}
\usepackage[normalem]{ulem}
\usepackage{capt-of}
\usepackage{hyperref}
\usepackage{minted}
\usepackage{fontspec}
\setmainfont{UT Sans}[
BoldFont={UT Sans Bold},
UprightFont={UT Sans Regular},
FontFace={m}{n}{UT Sans Medium}
]
% Fallback fake italic using a slant
\usepackage{etoolbox}
\newcommand{\fakeitalic}[1]{{\addfontfeatures{FakeSlant=0.18}#1}}
\AtBeginDocument{\renewcommand{\emph}[1]{\fakeitalic{#1}}}
\usepackage{minted}
\usepackage[a4paper,margin=3cm]{geometry}
\usepackage{lastpage}
\usepackage{fancyhdr}
\usepackage{lastpage}
\pagestyle{fancy}                % set fancy style
\fancyhf{}                        % clear all headers and footers
\cfoot{\thepage/\pageref*{LastPage}}  % center footer
\renewcommand{\headrulewidth}{0pt}    % remove header line
\renewcommand{\footrulewidth}{0pt}    % remove footer line
% also override plain page style for first page
\fancypagestyle{plain}{%
\fancyhf{}%
\cfoot{\thepage/\pageref*{LastPage}}%
\renewcommand{\headrulewidth}{0pt}%
\renewcommand{\footrulewidth}{0pt}%
}
\usepackage{url}
\author{Elias-Robert BAJCSI}
\date{}
\title{Laborator 5\\\medskip
\large PLFP}
\hypersetup{
 pdfauthor={Elias-Robert BAJCSI},
 pdftitle={Laborator 5},
 pdfkeywords={},
 pdfsubject={},
 pdfcreator={},
 pdflang={English}}
\begin{document}

\maketitle
\section{\texttt{./Cargo.toml}}
\label{sec:org1d891a5}
\begin{minted}[breaklines=true,breakanywhere=true,breaksymbol=,fontsize=\small,linenos=false]{toml}
[package]
name = "lab5"
version = "0.1.0"
edition = "2021"

[[bin]]
name = "ex1"
path = "src/bin/ex1.rs"
[[bin]]
name = "ex2"
path = "src/bin/ex2.rs"
[[bin]]
name = "ex3"
path = "src/bin/ex3.rs"
[[bin]]
name = "ex4"
path = "src/bin/ex4.rs"
\end{minted}
\section{\texttt{./src/bin/ex1.rs}}
\label{sec:org3f2eae8}
\begin{minted}[breaklines=true,breakanywhere=true,breaksymbol=,fontsize=\small,linenos=false]{rs}
use std::time::Instant;

fn pipeline(val: i32, pasi: &[fn(i32) -> i32]) -> i32 {
    pasi.iter().fold(val, |acc, f| f(acc))
}

fn pipeline_rev(val: i32, pasi: &[fn(i32) -> i32]) -> i32 {
    pasi.iter().rev().fold(val, |acc, f| f(acc))
}

fn adauga_1(x: i32) -> i32 { x + 1 }
fn dublu(x: i32) -> i32 { x * 2 }
fn la_patrat(x: i32) -> i32 { x * x }

fn main() {
    let a: &[fn(i32) -> i32] = &[adauga_1, dublu, la_patrat];
    let b: &[fn(i32) -> i32] = &[la_patrat, dublu, adauga_1];
    println!("{} == {}", pipeline_rev(5, a), pipeline(5, b));

    let t1 = Instant::now();
    let v: Vec<i32> = (1..=1000).filter(|n| n % 10 == 0).collect();
    let s1: i32 = v.iter().sum();
    let d1 = t1.elapsed();

    let t2 = Instant::now();
    let mut v2 = Vec::new();
    for n in 1..=1000 {
        if n % 10 == 0 {
            v2.push(n);
        }
    }
    let s2: i32 = v2.iter().sum();
    let d2 = t2.elapsed();

    println!("pipeline: {} {:?}", s1, d1);
    println!("for: {} {:?}", s2, d2);
}
\end{minted}
\section{\texttt{./src/bin/ex2.rs}}
\label{sec:orgf73a77f}
\begin{minted}[breaklines=true,breakanywhere=true,breaksymbol=,fontsize=\small,linenos=false]{rs}
fn compune_vec(fns: Vec<Box<dyn Fn(i32) -> i32>>) -> impl Fn(i32) -> i32 {
    move |x| fns.iter().fold(x, |acc, f| f(acc))
}

fn main() {
    let f = compune_vec(vec![
        Box::new(|x| x + 5),
        Box::new(|x| x * 3),
        Box::new(|x| x - 4),
    ]);

    println!("{}", f(10));
}
\end{minted}
\section{\texttt{./src/bin/ex3.rs}}
\label{sec:org5d2e2a6}
\begin{minted}[breaklines=true,breakanywhere=true,breaksymbol=,fontsize=\small,linenos=false]{rs}
fn cel_mai_bun_curry(compara: fn(i32, i32) -> bool) -> impl Fn(&[i32]) -> Option<i32> {
    move |nums: &[i32]| {
        nums.iter().copied().reduce(|acc, x| if compara(x, acc) { x } else { acc })
    }
}

fn este_mai_mic(a: i32, b: i32) -> bool { a < b }
fn este_mai_mare(a: i32, b: i32) -> bool { a > b }

fn suma_pare_pana_la(n: usize) -> i32 {
    (0..).filter(|x| x % 2 == 0).take(n).sum()
}

fn suma_pare_curry(n: usize) -> impl Fn() -> i32 {
    move || suma_pare_pana_la(n)
}

#[derive(Debug)]
struct Proiect {
    nume: String,
    activ: bool,
}

fn filtreaza_dupa_activ(activ: bool) -> impl Fn(&Proiect) -> bool {
    move |p| p.activ == activ
}

fn main() {
    let minim_partial = cel_mai_bun_curry(este_mai_mic);
    let maxim_partial = cel_mai_bun_curry(este_mai_mare);

    let v = [7, 3, 10, -1, 8];
    println!("{:?} {:?}", minim_partial(&v), maxim_partial(&v));

    println!("{} {}", suma_pare_pana_la(10), suma_pare_curry(10)());

    let proiecte = [
        Proiect { nume: "Licenta".into(), activ: true },
        Proiect { nume: "Restante".into(), activ: false },
        Proiect { nume: "Orar".into(), activ: true },
    ];

    let filtreaza_active = filtreaza_dupa_activ(true);
    let filtreaza_inactive = filtreaza_dupa_activ(false);

    let active: Vec<&str> = proiecte.iter().filter(|p| filtreaza_active(p)).map(|p| p.nume.as_str()).collect();
    let inactive: Vec<&str> = proiecte.iter().filter(|p| filtreaza_inactive(p)).map(|p| p.nume.as_str()).collect();

    println!("active: {:?}", active);
    println!("inactive: {:?}", inactive);
}
\end{minted}
\section{\texttt{./src/bin/ex4.rs}}
\label{sec:orgbc147fd}
\begin{minted}[breaklines=true,breakanywhere=true,breaksymbol=,fontsize=\small,linenos=false]{rs}
use std::cmp::Ordering;
use std::time::Instant;

#[derive(Debug, Clone)]
struct Proiect {
    nume: String,
    activ: bool,
    data_creare: u64,
    data_modificare: u64,
}

#[derive(Debug)]
struct Info {
    nume: String,
    detalii: String,
}

fn sorteaza_dupa(camp: fn(&Proiect) -> u64) -> impl Fn(&Proiect, &Proiect) -> Ordering {
    move |a, b| camp(a).cmp(&camp(b))
}

fn filtreaza_activ(activ: bool) -> impl Fn(&Proiect) -> bool {
    move |p| p.activ == activ
}

fn extrage_info(p: &Proiect) -> Info {
    Info { nume: p.nume.clone(), detalii: format!("modificat: {}", p.data_modificare) }
}

fn main() {
    let proiecte = vec![
        Proiect { nume: "Licenta".into(), activ: true, data_creare: 20240201, data_modificare: 20260508 },
        Proiect { nume: "Restante".into(), activ: false, data_creare: 20240115, data_modificare: 20240310 },
        Proiect { nume: "Orar".into(), activ: true, data_creare: 20231001, data_modificare: 20260224 },
        Proiect { nume: "Factura camin".into(), activ: true, data_creare: 20241105, data_modificare: 20260430 },
        Proiect { nume: "Adeverinta".into(), activ: false, data_creare: 20230920, data_modificare: 20231212 },
        Proiect { nume: "Dosar Erasmus".into(), activ: true, data_creare: 20250110, data_modificare: 20260502 },
        Proiect { nume: "Practica".into(), activ: true, data_creare: 20250303, data_modificare: 20260418 },
        Proiect { nume: "Bursa".into(), activ: false, data_creare: 20241010, data_modificare: 20250107 },
    ];

    let sort_creare = sorteaza_dupa(|p| p.data_creare);
    let sort_modificare = sorteaza_dupa(|p| p.data_modificare);

    let t1 = Instant::now();
    let mut p1: Vec<Proiect> = proiecte.iter().filter(|p| filtreaza_activ(true)(p)).cloned().collect();
    p1.sort_by(sort_modificare);
    let i1: Vec<Info> = p1.iter().map(extrage_info).collect();
    let d1 = t1.elapsed();

    let t2 = Instant::now();
    let mut p2 = Vec::new();
    for p in &proiecte {
        if p.activ {
            p2.push(p.clone());
        }
    }
    p2.sort_by_key(|p| p.data_modificare);
    let mut i2 = Vec::new();
    for p in &p2 {
        i2.push(extrage_info(p));
    }
    let d2 = t2.elapsed();

    for i in &i1 { println!("{}: {}", i.nume, i.detalii); }
    let mut dupa_creare = proiecte.clone();
    dupa_creare.sort_by(sort_creare);
    println!("primul creat: {}", dupa_creare[0].nume);
    println!("pipeline: {:?}", d1);
    println!("imperativ: {:?}", d2);
    println!("{}", i2.len());
}
\end{minted}
\end{document}
